\begin{python0}
from solutions import *; clear()
\end{python0}
\sectionthree{Operator as a Method}

\begin{consolethree}[escapeinside=||]
// Fraction.h
class Fraction
{
public:
    Fraction(int n=0, int d=1)
        : n_(n), d_(d)
    {}
    Fraction operator+(const Fraction &) const;
private:
    int n_, d_; // numerator,denominator
};
\end{consolethree}

\begin{consolethree}[escapeinside=||]
// Fraction.cpp
#include "Fraction.h"

Fraction Fraction::operator+(const Fraction & b) const
{
    return Fraction(n_ * b.d_ + d_ * b.n_, d_ * b.d_);
}
\end{consolethree}

Re-defining \textit{+}, so that we have \textit{+} for \textit{int},
\textit{double}, and \textit{Fraction}.

\begin{consolethree}[escapeinside=||]
#include "Fraction.h"

int main()
{
    Fraction a(1,3), b(1,4);
    Fraction c = a + b;
    Fraction d = a.operator+(b);
    return 0;
}
\end{consolethree}

Most programming languages allow you to define operators (with their own
syntax). There are languages that do not have operator overloading.
Example: Java.

\begin{consolethree}[escapeinside=||]
// Fraction.cpp
#include "Fraction.h"

Fraction Fraction::operator*(const Fraction & b) const
{
    return Fraction(n_ * b.n_, d_ * b.d_);
}
\end{consolethree}

Test it:

\begin{consolethree}[escapeinside=||]
#include "Fraction.h"

int main()
{
    Fraction a(1, 3), b(1, 4);
    Fraction c = a * b;
    return 0;
}
\end{consolethree}


